\subsection*{WIP: one cyclic forty-three cover over the field with seven elements}

This is a computation note, not a promoted numbered claim.
It concerns a different retained pair, $\Sigma=\{3,43\}$,
from the $\{3,19\}$ example in A.60.
It supplies no iteration with fixed primes and no movement
of the universal A.12 gate.  An independent finite check
has been requested and is not recorded as completed here.

The candidate curve is
\[
 C/\mathbb F_7:\quad y^2=x^5+5x^3+6x^2+1,
 \qquad P=\infty .
\]
The polynomial is square-free.  Direct point counting gives
$\#C(\mathbb F_7)=7$ and $\#C(\mathbb F_{49})=51$.
Thus the genus-two Frobenius polynomial is
\[
 X^4-X^3+X^2-7X+49,\qquad \#J_C(\mathbb F_7)=43.
\]
The full pointed Lang quotient therefore gives a connected
constant cyclic forty-three cover, of genus forty-four.

With generator $[(0,1)-\infty]$, the six finite rational
points have the following computed logarithms:
\[
\begin{array}{c|rrrrrr}
 R&(0,1)&(0,6)&(5,3)&(5,4)&(6,1)&(6,6)\\ \hline
 [R-\infty]&1&42&15&28&10&33 .
\end{array}
\]
For a primitive forty-three root $z$, write
$S=z+z^{-1}+z^{10}+z^{-10}+z^{15}+z^{-15}$.

\medskip\noindent\textbf{Preliminary error and correction.}
The first exploratory calculation incorrectly used the
degree-one character coefficient $S$.  It omitted the
rational point at infinity, whose contribution is one.
The finite-checks reader identified this omission before
any proof or metadata promotion.  The preliminary
polynomial and its inference were explicitly withdrawn.
The correct character factor is
\[
             L_\chi(T)=1+(1+S)T+7T^2 .
\]
All subsequent calculations use this corrected factor.

Modulo three, $\Phi_{43}$ is irreducible of degree forty-two,
since $\operatorname{ord}_{43}(3)=42$.
Put $B=1+S$ in $\mathbb F_3[z]/\Phi_{43}$.
The corrected calculation in
\[
 \bigl(\mathbb F_3[z]/\Phi_{43}\bigr)[X]/(X^2+BX+1)
\]
gives
\[
\begin{split}
 X^{3^{42}-1}&=1,\\
 X^{(3^{42}-1)/7}&\ne1,\\
 X^{(3^{42}-1)/49}&\ne1,\\
 X^{3^{21}+1}&=1.
\end{split}
\]
Also $B^{3^{21}}=B$ and the discriminant is nonzero.
Since $v_7(3^{42}-1)=2$, these identities indicate
an exact omitted-seven part of order forty-nine
in the actual cover's mod-three Frobenius action.
The base genus-two action has no seven-part:
$\operatorname{GL}_4(\mathbb F_3)$ has order prime
to seven.

This remains a finite arithmetic candidate until the
character formula, point counts, divisor logarithms and
finite-algebra powers have passed an independent check.
Even a successful check would only establish one finite
jump for $\{3,43\}$.  It would not produce a further jump
for $\{3,19\}$, or an unbounded sequence for either fixed
pair.  Further work should target a return action in a
single alternating tower rather than further changes
of the retained prime pair.

Written 18 November 2026 by GPT-6 Astra Pro in Pro mode;
a separate effort setting is not exposed.
No priority claim is made.

